\subsection{Dataset}
\label{sec:dataset}

\textbf{Documents.}
We collect 36 documents across three domains to serve as $K$: twelve Computer Science papers from arXiv, twelve legal opinions from U.S. federal appellate courts, and twelve medical case reports from PubMed Central.  To mitigate confounding factors from OLMo2's pre-training corpus, we select documents published near or after the cutoff date and verify their absence through the Infini-gram API~\citep{liu2024infini}.  Additional details on sourcing, preprocessing, and the corpus-membership check are provided in Appendix~\ref{Appendix:preprocess}.

\textbf{Views.} For our experiments, we syntheticaly generate auxiliary views with LLMs; inspired by our earlier observation and prior work~\citep{gunasekar2023textbooks, allen2024physics, jiang2024instruction}, we use textbooks, Stack Exchange--style Q\&A, and blogs as three distinct styles. Prerequisite knowledge is generated as textbook chapters that teach foundational concepts presupposed by the target document, explicitly instructed to avoid the source document's specific details. Contextual knowledge, however, is collected rather than generated: cited arXiv papers and cited judicial opinions; medical case reports lack an analogous citation structure, so we omit contextual knowledge for this domain. Paraphrases are generated with GPT-4.1, while auxiliary and prerequisite views are generated with GPT-5 and GPT-5-mini. More details on data collection, preprocessing, and prompts for synthetic data generation are provided in Appendix~\ref{Appendix:preprocess}.
\appendix
\section{Data}
\label{Appendix:preprocess}

\subsection{Dataset Construction}

\paragraph{Document collection.} The arXiv set was manually collected.  For medical documents, we use the NCBI E-utilities API to search PubMed Central for open-access case reports and download JATS XML.  For legal documents, we use the CourtListener REST API to search U.S. federal appellate opinions and download full opinion text.  Across domains, we filter candidate documents to obtain twelve target documents of suitable length, e.g., court opinions longer than 4000 tokens were rejected for efficient chunking.

\paragraph{Pre-training corpus check.} We check all 36 source documents against the Infini-gram API.  For each document, we query the title and ten randomly sampled body sentences against the OLMo2 pre-training mixture index, \texttt{v4\_olmo-mix-1124\_llama} and found zero counts.

\paragraph{Cleaning.}
For arXiv papers, we clean the raw LaTeX by removing comments, figures, tables, presentation-only commands, unresolved includes, page breaks, and bibliography/appendix material; expanding simple macros; extracting title, abstract, and main body; resolving revision commands; and repairing erroneous line breaks while preserving LaTeX environments.  Medical case reports are converted from JATS XML into sectioned plain text.  Legal opinions require additional PDF/OCR cleanup: we remove page headers, docket/page artifacts, decorative separators, extracted footnotes, and merge dangling sentences into paragraphs.

\paragraph{Contextual views.}
In our study, contextual views are cited works associated with the target document.  For arXiv documents, we resolve each target paper to an arXiv identifier, retrieve reference metadata using Semantic Scholar and OpenAlex, identify references with arXiv identifiers, download their source packages, and clean them with the same arXiv cleaning procedure.  For legal documents, we use CourtListener citation links to collect cited judicial opinions and apply the same legal-opinion cleaning procedure.  We do not construct contextual views for the medical documents because the case reports do not have cited works.

\subsection{Synthetic Data Generation}

\paragraph{Data statistics.}
Table~\ref{tab:dataset_artifact_inventory} summarizes the artifacts used in our experiments.  We generate 49 paraphrases per source document, yielding 1764 paraphrases in total.  For auxiliary views, $N$ counts the generated units obtained by splitting each view family into per-post blogs, per-question Stack Exchange Q\&A, and per-chapter textbooks. 

\begin{table}[t]
\centering
\scriptsize
\setlength{\tabcolsep}{3pt}
\begin{tabular}{@{}lrrr@{}}
\toprule
\textbf{Material} & \textbf{$N$} & \textbf{Avg. Len.} & \textbf{Total Tokens} \\
\midrule
\multicolumn{4}{@{}l}{\textit{Overall}} \\
Papers & 36 & 5858 & 210{,}888 \\
Paraphrases & 1764 & 6180 & 10{,}904{,}160 \\
Blogs & 225 & 1901 & 427{,}632 \\
Stack Exchange Q\&A & 453 & 1290 & 584{,}392 \\
Textbook Chapters & 282 & 2239 & 631{,}403 \\
Prerequisite Chapters & 582 & 4282 & 2{,}492{,}392 \\
Contextual Documents & 639 & 10416 & 6{,}655{,}888 \\
\midrule
\multicolumn{4}{@{}l}{\textit{Computer Science}} \\
Papers & 12 & 10550 & 126{,}598 \\
Paraphrases & 638 & 11002 & 7{,}019{,}384 \\
Blogs & 105 & 1920 & 201{,}648 \\
Stack Exchange Q\&A & 242 & 1460 & 353{,}321 \\
Textbook Chapters & 158 & 2872 & 453{,}833 \\
Prerequisite Chapters & 187 & 4555 & 851{,}724 \\
Contextual Documents & 506 & 10791 & 5{,}460{,}148 \\
\midrule
\multicolumn{4}{@{}l}{\textit{Medical}} \\
Papers & 12 & 3765 & 45{,}181 \\
Paraphrases & 588 & 3858 & 2{,}268{,}311 \\
Blogs & 72 & 1979 & 142{,}500 \\
Stack Exchange Q\&A & 117 & 1068 & 124{,}973 \\
Textbook Chapters & 61 & 1404 & 85{,}666 \\
Prerequisite Chapters & 220 & 4589 & 1{,}009{,}582 \\
Contextual Documents & 0 & -- & -- \\
\midrule
\multicolumn{4}{@{}l}{\textit{Legal}} \\
Papers & 12 & 3259 & 39{,}109 \\
Paraphrases & 589 & 3274 & 1{,}928{,}141 \\
Blogs & 48 & 1739 & 83{,}484 \\
Stack Exchange Q\&A & 94 & 1129 & 106{,}098 \\
Textbook Chapters & 63 & 1459 & 91{,}904 \\
Prerequisite Chapters & 175 & 3606 & 631{,}086 \\
Contextual Documents & 133 & 8991 & 1{,}195{,}740 \\
\bottomrule
\end{tabular}
\caption{Dataset statistics by material type and domain. $N$ counts text units; average lengths and totals are computed with the OLMo-2 tokenizer and rounded to the nearest token.}
\label{tab:dataset_artifact_inventory}
\end{table}

\paragraph{Paraphrases.}
Paraphrases are generated from cleaned source documents on a paragraph-level basis.  We preserve section headers verbatim and avoid paraphrasing LaTeX-only paragraphs. Domain-specific prompts are used for academic, legal, and medical text.  The academic prompt preserves LaTeX formatting, equations, proper nouns, titles, section headers, and technical terminology; the legal prompt preserves legal meaning, party names, citations, quoted language, dates, docket numbers, and legal terms of art; and the medical prompt preserves diagnoses, medications, doses, routes, timelines, lab values, imaging findings, procedures, outcomes, and clinical terminology. We use GPT-4.1 to generate paraphrases with a temperature of 1 and a top-p of 0.975.

\paragraph{Auxiliary views.}
For each document, we generate synthetic auxiliary views conditioned on the target document and intended to restate, explain, or pedagogically reorganize the same document-level content.  The Computer Science setting generates textbook chapters for college students with basic machine-learning background, Stack Exchange--style questions from a confused student followed by grounded answers, and technical blog posts for a broader technical audience.  The medical setting adapts these formats to clinical education: case-based textbook sections for medical students or residents, clinical teaching Q\&A, and clinical blog posts.  The legal setting uses casebook/treatise-style chapters for law students, Law Stack Exchange-style Q\&A, and analytic legal commentary.  For each view family, we first generate an outline, list of questions, or list of blog ideas, and then generate the individual chapters, answers, or posts from that plan. We use GPT-5 to generate the outlines and GPT-5-mini to generate the contents of the chapters, questions, or posts.

\paragraph{Prerequisite views.}
Prerequisite views are generated separately from auxiliary views.  For arXiv papers, we use a curriculum-design prompt to create prerequisite textbook chapters that teach the foundations needed to understand the paper while excluding the paper's own novel ideas.  For medical case reports, the prompt asks for general medical background--for example relevant anatomy, pathophysiology, pharmacology, diagnostic interpretation, differential diagnosis, and standard management--while explicitly forbidding patient-specific chronology, workup, treatment course, outcome, or novel observations from the source case.  For legal opinions, the prompt similarly asks for doctrinal and procedural background while excluding the source case's parties, facts, outcome, holding, and novel reasoning. When useful for legal background, we additionally include landmark or doctrinally foundational cited opinions as context for the generated prerequisite chapters. We use GPT-5 to generate the outlines and GPT-5-mini to generate the contents of the chapters.
